% =====================================================================
%  Appendix A4 — Bibliography
% =====================================================================
\chapter*{Bibliography}
\addcontentsline{toc}{chapter}{Bibliography}
\markboth{Bibliography}{Bibliography}

\noindent
\small
This bibliography has three parts. Entries \textbf{[1]--[28]} are reproduced
from the bibliography of \texttt{proofs/claim\_a\_3d\_proof\_attempt.tex}, in
the document's own order and wording; those are the works the programme's
mathematics actually cites, and the five keys this book cites directly ---
\cite{KNSS2009}, \cite{SereginSverak2009}, \cite{Seregin2012},
\cite{LadyzhenskayaSeregin1999}, \cite{AlbrittonBarker2019} --- are among
them. Entries \textbf{[29]--[36]} are classical works this book refers to by
name in the prose without the proof document citing them formally. Entry
\textbf{[37]} is a group entry for the contemporaneous forced-blow-up results
of Chapter~\ref{ch:openai}; it is deliberately descriptive rather than
precise, for the reason given there --- those claims were read through
secondary summaries, not preprints, and this book does not pretend to a
citation it cannot verify.

\medskip
\noindent
Two of the document's own citations are \emph{miscited} at the point of use,
and that is recorded in \texttt{rem:s48-miscitations} and in
Chapter~\ref{ch:cz-log}, \S\ref{sec:cz-miscitations}. The entries below are
correct; it is the sentences citing them, at lines 6798, 6810, 6848 and 7689
of the proof document, that are not.

\begin{thebibliography}{99}

\bibitem{CKN1982}
  L.~Caffarelli, R.~Kohn, L.~Nirenberg,
  Partial regularity of suitable weak solutions of the Navier--Stokes
  equations,
  \emph{Comm. Pure Appl. Math.}, 35(6):771--831, 1982.

\bibitem{BenichouBrezis}
  Ph.~B\'enilan, H.~Brezis, M.~G.~Crandall,
  A semilinear equation in \(L^1(\R^N)\),
  \emph{Ann. Scuola Norm. Sup. Pisa}, 2(4):523--555, 1975.

\bibitem{Fefferman2000}
  C.~L.~Fefferman,
  \emph{Existence and Smoothness of the Navier--Stokes Equation},
  Clay Mathematics Institute, 2000.
  \textsf{[The official problem statement. Its four alternatives (A)--(D) are
  the subject of Chapter~\ref{ch:openai} and of
  Chapter~\ref{ch:where-it-stands}, \S\ref{sec:naming}.]}

\bibitem{Giaquinta1982}
  M.~Giaquinta, E.~Giusti,
  On the regularity of the minima of variational integrals,
  \emph{Acta Math.}, 148:31--46, 1982.

\bibitem{Lin1998}
  F.~Lin,
  A new proof of the Caffarelli--Kohn--Nirenberg theorem,
  \emph{Comm. Pure Appl. Math.}, 51(3):241--257, 1998.

\bibitem{Prodi1959}
  G.~Prodi,
  Un teorema di unicit\`a per le equazioni di Navier--Stokes,
  \emph{Ann. Mat. Pura Appl.}, 48:173--182, 1959.

\bibitem{Serrin1962}
  J.~Serrin,
  On the interior regularity of weak solutions of the Navier--Stokes
  equations,
  \emph{Arch. Rat. Mech. Anal.}, 9:187--195, 1962.
  \textsf{[The strict subcritical range \(2/s'+3/s<1\) proved here is what
  decides hypothesis \textbf{(R)} at \Sn{48}. See
  Chapter~\ref{ch:cz-log}, \S\ref{sec:s48-R}.]}

\bibitem{Stein1970}
  E.~M.~Stein,
  \emph{Singular Integrals and Differentiability Properties of Functions},
  Princeton Univ. Press, 1970.
  \textsf{[The \(L^\infty\to\mathrm{BMO}\) endpoint of a Calder\'on--Zygmund
  operator; the reason Wall~1 exists.]}

\bibitem{Stredulinsky1984}
  E.~W.~Stredulinsky,
  Higher integrability from reverse H\"older inequalities,
  \emph{Indiana Univ. Math. J.}, 29:407--413, 1980.

\bibitem{Krylov1996}
  N.~V.~Krylov,
  \emph{Lectures on Elliptic and Parabolic Equations in H\"older Spaces},
  Graduate Studies in Mathematics, Amer. Math. Soc., 1996.

\bibitem{BenilBrezis1975}
  Ph.~B\'enilan, H.~Br\'ezis, M.~G.~Crandall,
  A semilinear equation in \(L^1(\R^N)\),
  \emph{Ann. Sc. Norm. Sup. Pisa}, 2(4):523--555, 1975.
  \textsf{[A duplicate of [2] under a second key; both are in the proof
  document's bibliography and both are reproduced here rather than silently
  merged.]}

\bibitem{ConstantinFefferman1993}
  P.~Constantin, C.~Fefferman,
  Direction of vorticity and the problem of global regularity for the
  Navier--Stokes equations,
  \emph{Indiana Univ. Math. J.}, 42(3):775--789, 1993.
  \textsf{[The origin of the coherence criterion the whole of
  Part~IV is built on, and the source of the calibration
  \(\sstar\sim1\).]}

\bibitem{Constantin1994}
  P.~Constantin,
  Geometric statistics in turbulence,
  \emph{SIAM Review}, 36(1):73--98, 1994.
  \textsf{[The direction equation of \texttt{lem:direction-eq} and the exact
  kernel cancellation used in \texttt{prop:localized-cf}.]}

\bibitem{KozonoTaniuchi2000}
  H.~Kozono, Y.~Taniuchi,
  Bilinear estimates in BMO and the Navier--Stokes equations,
  \emph{Math. Z.}, 235(1):173--194, 2000.

\bibitem{NecasRuzickaSverak1996}
  J.~Ne\v{c}as, M.~R\r{u}\v{z}i\v{c}ka, V.~\v{S}ver\'ak,
  On Leray's self-similar solutions of the Navier--Stokes equations,
  \emph{Acta Math.}, 176:283--294, 1996.

\bibitem{EscauriazaSereginSverak2003}
  L.~Escauriaza, G.~Seregin, V.~\v{S}ver\'ak,
  \(L_{3,\infty}\)-solutions of Navier--Stokes equations and backward
  uniqueness,
  \emph{Uspekhi Mat. Nauk}, 58:3--44, 2003.

\bibitem{KNSS2009}
  G.~Koch, N.~Nadirashvili, G.~Seregin, V.~\v{S}ver\'ak,
  Liouville theorems for the Navier--Stokes equations and applications,
  \emph{Acta Math.}, 203(1):83--105, 2009.
  \textsf{[Two distinct uses: the two-dimensional bounded-ancient Liouville
  theorem that closes the aligned horn, and the \(L^\infty\)-based local
  theory of \S4 that supplies the quantitative form of \texttt{thm:s48-R}.]}

\bibitem{Seregin2012}
  G.~Seregin,
  A certain necessary condition of potential blow up for Navier--Stokes
  equations,
  \emph{Comm. Math. Phys.}, 312(3):833--845, 2012.
  \textsf{[\textbf{Miscited} at lines 6810, 6848 and 7689 of the proof
  document as authority for Type-I local regularity technology. It is an
  \(L^3\) critical-endpoint blow-up result and contains no such theorem.
  \textbf{(N)} is unaffected: it remains assumed.]}

\bibitem{LadyzhenskayaSeregin1999}
  O.~A.~Ladyzhenskaya, G.~A.~Seregin,
  On partial regularity of suitable weak solutions to the three-dimensional
  Navier--Stokes equations,
  \emph{J. Math. Fluid Mech.}, 1(4):356--387, 1999.
  \textsf{[\textbf{Loosely attributed} at line~6798 as a ``sharpening by
  local-pressure theory''. It is an \(\varepsilon\)-regularity paper for
  suitable weak solutions requiring \(p\in L^{3/2}_{\mathrm{loc}}\) and a
  smallness hypothesis. It is not the tool that discharges \textbf{(R)}.]}

\bibitem{Struwe1988}
  M.~Struwe,
  On partial regularity results for the Navier--Stokes equations,
  \emph{Comm. Pure Appl. Math.}, 41(4):437--458, 1988.

\bibitem{Takahashi1990}
  S.~Takahashi,
  On interior regularity criteria for weak solutions of the Navier--Stokes
  equations,
  \emph{Manuscripta Math.}, 69(3):237--254, 1990.
  \textsf{[With [20], the borderline equality case of Serrin's criterion ---
  which \Sn{48} explicitly does \emph{not} need, because the available
  exponent pair is strictly subcritical.]}

\bibitem{SereginSverak2009}
  G.~Seregin, V.~\v{S}ver\'ak,
  On type~I singularities of the local axi-symmetric solutions of the
  Navier--Stokes equations,
  \emph{Comm. Partial Differential Equations}, 34(2):171--201, 2009.
  (arXiv:0804.1803; Definition~2.3, Theorem~2.4, Theorem~2.8 and Remark~2.9
  are the statements used in \texttt{sec:hypR-s48}.)

\bibitem{AlbrittonBarker2019}
  D.~Albritton, T.~Barker,
  On local type~I singularities of the Navier--Stokes equations and Liouville
  theorems,
  \emph{J. Math. Fluid Mech.}, 21(3):Art.~43, 2019. (arXiv:1811.00502.)

\bibitem{ChenGaldiPoggiSchikorra2026}
  R.~M.~Chen, G.~P.~Galdi, B.~Poggi, A.~Schikorra,
  On Serrin interior regularity criterion for Navier--Stokes equations,
  arXiv:2606.24733, 2026.
  \textsf{[Quoted in \texttt{sec:hypR-s48} for the exact statement of
  Serrin's criterion and its attribution.]}

\bibitem{Kwon2021}
  H.~Kwon,
  The role of the pressure in the regularity theory for the Navier--Stokes
  equations,
  arXiv:2104.03160, 2021.

\bibitem{Giusti2003}
  E.~Giusti,
  \emph{Direct Methods in the Calculus of Variations},
  World Scientific, 2003. (Iteration Lemma~6.1; hole-filling.)
  \textsf{[The linear iteration lemma that \texttt{prop:s41-nocontraction}
  proves unavailable: it needs a contraction factor below one.]}

\bibitem{LadyzhenskayaUraltseva1968}
  O.~A.~Ladyzhenskaya, N.~N.~Ural'tseva,
  \emph{Linear and Quasilinear Elliptic Equations},
  Academic Press, 1968. (Chapter~II: De~Giorgi level truncation.)
  \textsf{[The scheme \texttt{rem:level-iteration} invoked by name, and that
  Chapter~\ref{ch:annulus} shows does not apply, because the truncation here
  is in an auxiliary variable on a fixed band and no set shrinks.]}

% --- classical works referred to by name in this book ---

\bibitem{Leray1934}
  J.~Leray,
  Sur le mouvement d'un liquide visqueux emplissant l'espace,
  \emph{Acta Math.}, 63:193--248, 1934.
  \textsf{[Global existence of weak solutions. The programme's own answer to
  ``does this prove existence?'' is \emph{no}: existence is here, and the open
  question is smoothness. Recorded in the proof document at \Sn{35}.]}

\bibitem{BKM1984}
  J.~T.~Beale, T.~Kato, A.~Majda,
  Remarks on the breakdown of smooth solutions for the 3-D Euler equations,
  \emph{Comm. Math. Phys.}, 94(1):61--66, 1984.
  \textsf{[The BKM continuation criterion. Defect~D5 of the \Sn{31} audit is
  that \texttt{thm:sigma-gronwall} assumes it, hence assumes an equivalent of
  its own conclusion.]}

\bibitem{FujitaKato1964}
  H.~Fujita, T.~Kato,
  On the Navier--Stokes initial value problem I,
  \emph{Arch. Rat. Mech. Anal.}, 16:269--315, 1964.
  \textsf{[Local well-posedness. Used twice: for the near-shear perturbation
  theorem of \Sn{21}, and as the engine of
  \texttt{prop:regime-A-impossible}.]}

\bibitem{Gehring1973}
  F.~W.~Gehring,
  The \(L^p\)-integrability of the partial derivatives of a quasiconformal
  mapping,
  \emph{Acta Math.}, 130:265--277, 1973.
  \textsf{[Gehring's lemma: reverse H\"older gives \(\delta_0>0\). Step~5 of
  Route~B.]}

\bibitem{LSU1968}
  O.~A.~Ladyzhenskaya, V.~A.~Solonnikov, N.~N.~Ural'tseva,
  \emph{Linear and Quasilinear Equations of Parabolic Type},
  Amer. Math. Soc., 1968.
  \textsf{[The parabolic embedding \(L^\infty L^2\cap L^2H^1\hookrightarrow
  L^{10/3}(Q)\) used by \texttt{prop:second-rung}.]}

\bibitem{Stampacchia1966}
  G.~Stampacchia,
  \emph{\'Equations elliptiques du second ordre \`a coefficients
  discontinus},
  Presses de l'Universit\'e de Montr\'eal, 1966.
  \textsf{[The nonlinear iteration lemma that
  \texttt{rem:s41-whichlemma} shows is equally unavailable: it needs a
  superlinear gain and small initial data, and the programme has neither.]}

\bibitem{Hamilton1993}
  R.~S.~Hamilton,
  A matrix Harnack estimate for the heat equation,
  \emph{Comm. Anal. Geom.}, 1(1):113--126, 1993.
  \textsf{[``Hamilton's trick'': at an interior spatial maximum
  \(\Delta\abs{\omega}\le0\). The step that turns the magnitude equation into
  the log-time ODE of \texttt{prop:logtime-necessary}.]}

\bibitem{CCRT2007}
  S.~I.~Chernyshenko, P.~Constantin, J.~C.~Robinson, E.~S.~Titi,
  A posteriori regularity of the three-dimensional Navier--Stokes equations
  from numerical computations,
  \emph{J. Math. Phys.}, 48(6):065204, 2007.
  \textsf{[The shape of the \Sn{39} a posteriori diagnostic. The module's own
  documentation is emphatic that it does \emph{not} produce a proof: a
  rigorous computer-assisted argument needs validated interval arithmetic on
  the residual and on every norm entering the Gr\"onwall estimate, and the
  implementation is ordinary floating point.]}

\bibitem{WangEtAl}
  Y.~Wang and co-authors, self-similar blow-up profiles for the
  Cordoba--Cordoba--Fontelos and Boussinesq models.
  \textsf{[The source of the \(\lambda\) values the M3 sweep uses:
  CCF stable \(1.1808\), CCF 1st unstable \(0.6057\), CCF 2nd unstable
  \(0.4703\); Boussinesq stable \(1.9206\), 1st \(1.3991\), 3rd \(1.1843\).
  Cited in the programme's own specification, not in the proof document;
  the bibliographic details are given as the specification records them and
  have not been independently verified for this book.]}

\bibitem{ForcedBlowup2026}
  \textsf{[Group entry, deliberately descriptive.]}
  Two contemporaneous efforts on finite-time blow-up for \emph{forced}
  equations, discussed in Chapter~\ref{ch:openai}. \emph{(i)} L.~Alp\"oge and
  T.~Buckmaster, extending a programme of A.~C\'ordoba and
  L.~Mart\'inez-Zoroa from rough to smooth forcing: released results for the
  incompressible porous medium equation, for Boussinesq, and for
  three-dimensional incompressible Euler, formalised in Lean, with a further
  hypo-dissipative Navier--Stokes claim stated as believed but unreleased.
  \emph{(ii)} A claim by OpenAI for the full three-dimensional incompressible
  Navier--Stokes equations at arbitrary \(\nu>0\), from rest, with bounded
  kinetic energy and unbounded velocity in finite time, with a Lean
  formalisation, presented on their own pages as establishing Fefferman's
  statement (C) and also (D).
  \textsf{[No preprint identifiers are given because this book read these
  claims through secondary summaries, and Chapter~\ref{ch:openai} says so.
  The book reports the claims and does not adjudicate them.]}

\end{thebibliography}
